\section{Quantum toroidal \texorpdfstring{$\mathfrak{gl}_1$}{gl1},
K3 equivariance, and the \texorpdfstring{$N=2$}{N=2} root}
\label{sec:qtgl1-k3-n2-root}

\providecommand{\DIM}{\mathrm{DIM}}
\providecommand{\gghone}{\widehat{\widehat{\mathfrak{gl}}}_1}
\providecommand{\Uqtgltor}{U_{q_1,q_2,q_3}(\gghone)}
\providecommand{\CoHA}{\mathrm{CoHA}}
\providecommand{\Irr}{\mathrm{Irr}}

The quantum-toroidal input is the Ding-Iohara-Miki algebra on the
Calabi-Yau parameter surface
\[
  \mathcal L_{\mathrm{Miki}}
  =
  \{(q_1,q_2,q_3)\in(\mathbb C^\times)^3:
  q_1q_2q_3=1\}.
\]
The equation cuts out a two-dimensional parameter surface.  Evaluations
at distinguished points are additional data on this surface.  The toric model
$\mathbb C^3$ supplies three equivariant weights
$\epsilon_i=\log q_i$ with
$\epsilon_1+\epsilon_2+\epsilon_3=0$.  Miki's symmetry belongs to this
toric algebraic situation, where the three weights are present before
specialisation.

\subsection{Parameter strata}
\label{subsec:qtet-parameter-strata}

\begin{definition}[Miki surface and self-dual slice]
The \emph{Miki surface} is
\[
 \mathcal L_{\mathrm{Miki}}
 = \operatorname{Spec}
 \mathbb C[q_1^{\pm1},q_2^{\pm1},q_3^{\pm1}]
 /(q_1q_2q_3-1).
\]
The self-dual slice is the one-parameter sublocus
\[
  q_1=q,\qquad q_2=q^{-1},\qquad q_3=1.
\]
The point used at level \(N=2\) is
\[
  (q_1,q_2,q_3)=(-1,-1,1).
\]
\end{definition}

\begin{proposition}[Generic symmetry and specialisation]
\label{prop:qtet-miki-specialisation}
For the rank-one toric algebra \(\Uqtgltor\), the current
presentation is symmetric in the three Calabi-Yau weights before a
current direction is singled out.  At the point
\((q_1,q_2,q_3)=(-1,-1,1)\), the transposition \(q_1\leftrightarrow
q_2\) preserves the displayed current realisation.  A cyclic
permutation sends the displayed current realisation to a different
choice of distinguished denominator and therefore requires an explicit
intertwiner.  The numerical equality \(q_1q_2q_3=1\) is insufficient.
\end{proposition}

\begin{proof}
Miki \(2007\) and the Feigin-Hashizume-Hoshino-Shiraishi-Yanagida
presentation define the algebra over the two-parameter surface
\(\mathcal L_{\mathrm{Miki}}\).  The usual current kernel
\[
  \omega_{\DIM}(x)
  =
  \frac{(1-q_1x)(1-q_2x)}{1-q_3x}
\]
distinguishes \(q_3\) in notation.  Swapping \(q_1\) and \(q_2\)
leaves this formula fixed.  Moving \(q_3\) into the numerator changes
the chosen current chart.  Thus a statement about an \(S_3\)-action at
a special point requires identified current charts in addition to the
product constraint.
\end{proof}

\begin{remark}[K3 equivariance]
\label{rem:qtet-k3-equivariance}
The K3 quantum-toroidal branch requires an external toric chart for a
\((\mathbb C^\times)^3\)-weight package.  For a generic K3 surface,
\(\operatorname{Aut}^{\circ}(K3)=\{1\}\).  For \(K3\times E\), the
connected automorphism group contributed by the elliptic factor is
\(E\).  Since \(E\) is complete, every morphism
\(\mathbb G_m\to E\) is constant.
Maulik-Okounkov stable-envelope constructions on K3 therefore require
an explicit reduced equivariance package, such as a Kummer or ADE
stratum, a formal equivariant replacement, or a specified elliptic
factor.  A quantum-toroidal statement for K3 that uses an ambient
three-dimensional torus must name that extra structure.
\end{remark}

\subsection{The positive half from \texorpdfstring{$\CoHA(\mathbb C^3)$}{CoHA(C3)}}
\label{subsec:qtet-coha-positive-half}

\begin{theorem}[Positive-half identification]
\label{thm:qtet-coha-positive-half}
Let \(T=(\mathbb C^\times)^3/\mathbb C^\times_{\mathrm{diag}}\) act on
\(\mathbb C^3\), and let
\(\epsilon_1+\epsilon_2+\epsilon_3=0\) be the induced Calabi-Yau
condition.  After localisation to the fraction field of the equivariant
parameters, the critical cohomological Hall algebra of the three-loop
Jordan quiver with potential is the shuffle algebra with kernel
\[
  \omega(z,w)
  =
  \frac{(z-w-\epsilon_1)(z-w-\epsilon_2)(z-w-\epsilon_3)}
       {(z-w)^3}.
\]
Schiffmann-Vasserot identify the localised shuffle algebra with the
positive shuffle algebra of the affine Yangian, Kapranov-Vasserot and
the critical-CoHA constructions supply the \(\mathbb C^3\) critical
cohomology model, and Tsymbaliuk identifies the shuffle
presentation with the Drinfeld-current presentation.  The resulting
algebra is the positive Drinfeld half
\[
  \CoHA(\mathbb C^3)\simeq Y^+_{\epsilon_1,\epsilon_2,\epsilon_3}
  (\gghone).
\]
The full affine Yangian is obtained by adjoining Cartan and negative
halves.  The vertex algebra \(\mathcal W_{1+\infty}\) appears through
an evaluation representation of the completed double.
\end{theorem}

\begin{proof}
Schiffmann-Vasserot identify the localised cohomological Hall shuffle
algebra with the positive shuffle algebra attached to the instanton
surface.  Tsymbaliuk identifies the corresponding shuffle algebra with
the positive Drinfeld half in the triangular decomposition
\[
  Y(\gghone)=Y^-\otimes Y^0\otimes Y^+ .
\]
The multiplication is the CoHA shuffle product.  The negative half,
Cartan half, Hopf pairing, and completed Drinfeld double are additional
data.
\end{proof}

\begin{corollary}[Positive half, double, and coproduct]
\label{cor:qtet-positive-half-no-hopf}
The Chevalley involution and Drinfeld coproduct used by the full
quantum-toroidal algebra are structures of a chosen Hopf completion or
of the Drinfeld double
\[
  D(Y^+)=Y^-\bowtie Y^0\bowtie Y^+ .
\]
For a Calabi-Yau input of dimension \(d\geq3\), the native
specialised chiral output is \(E_1\).  Any \(E_2\)-structure belongs
to a centre, double, or representation category after the relevant
hypotheses have been imposed.
\end{corollary}

\subsection{The \texorpdfstring{$N=2$}{N=2} current specialisation}
\label{subsec:qtet-n2-current}

\begin{proposition}[The current kernel at \texorpdfstring{$q=-1$}{q=-1}]
\label{prop:qtet-n2-current-kernel}
On the self-dual slice \(q_1=q\), \(q_2=q^{-1}\), \(q_3=1\), the
Ding-Iohara current kernel becomes
\[
  \omega_{\DIM}(x)
  =
  \frac{(1-qx)(1-q^{-1}x)}{1-x}.
\]
At \(q=-1\),
\[
  \omega_{\DIM}(x)=\frac{(1+x)^2}{1-x}.
\]
Thus the \(N=2\) point is singular in this current chart at \(x=1\).
\end{proposition}

\begin{proof}
Substitution gives the formula.  The pole at \(x=1\) is visible before
any representation-theoretic quotient is taken.
\end{proof}

\begin{remark}[Central rank and integral form]
\label{rem:qtet-central-rank-integral-form}
The pole of \(\omega_{\DIM}(x)\) at \(x=1\) records singular behaviour
of the current chart and the integral form at the root of unity.  The
central rank of the specialised algebra is determined only after the
specialised integral form, its Cartan completion, and the centre after
the root-of-unity quotient have been constructed.  The displayed kernel
proves exactly the singular-current statement.
\end{remark}

\subsection{K3 branches and the BKM boundary}
\label{subsec:qtet-k3-branches}

\begin{proposition}[Two K3 quantum-group branches]
\label{prop:qtet-two-k3-branches}
The Mukai stable-envelope branch and the \(K3\times E\) BKM branch are
different objects.
\begin{enumerate}[label=\textup{(\roman*)}]
\item The Mukai branch is the conjectural Mukai self-mirror K3 Yangian
or K3 quantum-toroidal algebra built from the Mukai lattice and
Maulik-Okounkov stable-envelope operators on Hilbert-scheme towers.
\item The BKM branch is the completed Hall-Drinfeld double attached
to \(K3\times E\), with Borcherds denominator data controlled by
\(\Delta_5\) and by the \(\mathfrak g_{\Delta_5}\) root
multiplicities.
\end{enumerate}
The BKM object is the Hall-Drinfeld/Borcherds object.  A map from a
positive K3 current block to the BKM boundary requires the Hall pairing,
negative half, completion, imaginary-root relations, and Igusa
orientation datum.
\end{proposition}

\begin{proof}
The Mukai branch is a stable-envelope construction on K3 Hilbert
schemes.  The BKM branch is a Hall-Drinfeld double for the
\(K3\times E\) Calabi-Yau threefold and its automorphic denominator.
One construction supplies positive current operators; the other
requires the doubled Hall algebra and the Borcherds root datum.  The
input categories, completions, and denominator data are different.
\end{proof}

\begin{proposition}[The four \(K3\times E\) invariants]
\label{prop:qtet-k3xe-four-invariants}
The \(K3\times E\) boundary uses four distinct invariants:
\[
  \kappa_{\mathrm{cat}}(K3\times E)=0,\qquad
  \kappa_{\mathrm{ch}}^{\mathrm{Heis}}=3,\qquad
  \kappa_{\mathrm{BKM}}(\Delta_5)=\frac{c_1(0)}2=\frac{10}{2}=5,
  \qquad
  \kappa_{\mathrm{fiber}}=24.
\]
The first is the K\"unneth product
\(\chi(\mathcal O_{K3})\chi(\mathcal O_E)=2\cdot0=0\).  The second is
the Heisenberg-Mukai specialisation.  The third is the Borcherds
weight of the Gritsenko-Nikulin denominator.  The fourth is the rank of
the Mukai lattice.
\end{proposition}

\begin{proposition}[Lie-level and vertex-level separation]
\label{prop:qtet-lie-vertex-separation}
At the abelian K3 Yangian stratum, the root Lie quotient is abelian:
root brackets factor through the central Heisenberg summand.  The
vertex and Hopf data carry extra structure.  The Heisenberg fields
\(J_a(z)\), \(a\in\Lambda_{\mathrm{Muk}}\), satisfy
\[
  J_a(z)J_b(w)\sim
  \frac{\langle a,b\rangle_{\mathrm{Muk}}}{(z-w)^2},
\]
and the generic spectral coproduct is controlled by the diagonal
Maulik-Okounkov factor
\[
  g_a(u)=\frac{u-h_a}{u+h_a}.
\]
Thus the Lie quotient, Heisenberg vertex algebra, and spectral Hopf
algebra are three separate levels of structure.
\end{proposition}

\begin{proposition}[Stable-envelope \(R\)-matrix data]
\label{prop:qtet-stable-envelope-rmatrix}
The Maulik-Okounkov \(R\)-matrix entering the K3 branch is determined
only after the following data are fixed:
\[
  (T,\mathfrak C,\mathsf{pol},\mathsf{slope},
  K_T(\mathrm{Hilb}^{\bullet}S)).
\]
Here \(T\) is an algebraic torus acting on the local surface chart,
\(\mathfrak C\) is a chamber, \(\mathsf{pol}\) is a polarization, and
\(\mathsf{slope}\) is the slope parameter.  The local operator is
\[
  R^{\mathrm{MO}}_{\mathfrak C_+,\mathfrak C_-}(u)
  =
  \bigl(\mathrm{Stab}^{T}_{\mathfrak C_-}(u)\bigr)^{-1}
  \mathrm{Stab}^{T}_{\mathfrak C_+}(u).
\]
For K3 this is available on ADE/Kummer charts or after a formal
reduced-equivariant replacement.  A global K3 or \(K3\times E\)
statement requires descent of these local \(R\)-matrices and their
cocycles.  Super-Yangian extensions, including Hodge-parity
\(\mathfrak{so}(4\mid20)\)-type lifts, carry conjectural status until
that descent and the parity form are constructed.
\end{proposition}

\begin{remark}[Elliptic and toroidal hypotheses]
\label{rem:qtet-elliptic-hypotheses}
An elliptic quantum-toroidal statement for \(K3\times E\) needs two
parameters named explicitly: the trigonometric Mukai parameter and the
elliptic modulus of \(E\).  The theta-kernel expression
\[
  G_{K3\times E}(x;q,\tau)
  =
  \prod_{a=1}^{24}
  \frac{\vartheta_1(xq_a;\tau)}
       {\vartheta_1(xq_a^{-1};\tau)}
\]
is meaningful only after the \(24\) Mukai parameters, their product
constraint, and the stable-envelope equivariance have been fixed.  It
is separate from any torus action on K3.
\end{remark}

\subsection{The coefficient \texorpdfstring{$324$}{324}}
\label{subsec:qtet-324}

\begin{proposition}[G\"ottsche and Fock counts]
\label{prop:qtet-324-gottsche-fock}
The coefficient of \(q^2\) in
\[
  \prod_{m\geq 1}(1-q^m)^{-24}
\]
is
\[
  24+\binom{24+1}{2}=24+300=324.
\]
Consequently
\[
  \chi(\mathrm{Hilb}^2(K3))=324
\]
by the G\"ottsche-Grojnowski-Nakajima Hilbert-scheme formula, and
the degree-\(2\) Fock space of \(24\) free Heisenberg generators has
dimension \(324\).
\end{proposition}

\begin{proof}
The term \(24q^2\) comes from choosing one \(q^2\)-mode from one of
the \(24\) factors.  The term \(\binom{24+1}{2}q^2\) comes from two
\(q\)-modes with repetition among \(24\) colours.  The same generating
series is the Hilbert-scheme Euler-characteristic series for K3.
\end{proof}

\begin{definition}[Four appearances of \(324\)]
\label{def:qtet-four-324s}
\[
\begin{array}{ll}
324_{\mathrm{Hilb}} &
  = \chi(\mathrm{Hilb}^2(K3)),\\[2pt]
324_{\mathrm{Fock}} &
  = \dim \mathrm{Fock}(H_{\mathrm{Muk}})_2,\\[2pt]
324_{\mathrm{block}} &
  = 5+44+44+231,\\[2pt]
324_{\mathrm{irr}} &
  = \left|\Irr(\mathcal C_{N=2})\right|
  \quad\text{for a specified root-of-unity tensor category }
  \mathcal C_{N=2}.
\end{array}
\]
The first two are equal by theorem.  The third is an arithmetical
decomposition until its sectors are constructed.  The fourth is a
representation-theoretic claim and depends on the chosen category
\(\mathcal C_{N=2}\).
\end{definition}

\begin{proposition}[Root-of-unity module count as proof obligation]
\label{prop:qtet-324-proof-obligation}
The equality
\[
  324_{\mathrm{irr}}=324_{\mathrm{Hilb}}
\]
requires more than the G\"ottsche coefficient.  It requires:
\begin{enumerate}[label=\textup{(\roman*)}]
\item a specified root-of-unity algebra, either the K3 positive
stable-envelope specialisation or the BKM small form;
\item an integral form and root-of-unity quotient;
\item a semisimplification or logarithmic replacement;
\item a classification of simple objects in that category;
\item a functor identifying the charge-\(2\) Fock basis with those
simple objects.
\end{enumerate}
Feigin-Tsymbaliuk finite-type and Fock-type module theory supplies
representation-theoretic input for quantum-toroidal
\(\mathfrak{gl}_1\).  A classification of all finite-dimensional
irreducibles for the K3 or BKM root-of-unity specialisation remains a
separate theorem.
\end{proposition}

\subsection{The \texorpdfstring{$S$}{S}-matrix}
\label{subsec:qtet-s-matrix}

\begin{proposition}[A Hadamard block count is full rank]
\label{prop:qtet-hadamard-full-rank}
A block diagonal \(324\times324\) matrix consisting of \(24\) ordinary
\(2\times2\) Hadamard blocks and \(276\) nonzero \(1\times1\) blocks
has rank
\[
  24\cdot 2+276=324.
\]
Therefore such a block decomposition proves full rank for that matrix.
\end{proposition}

\begin{proof}
An ordinary \(2\times2\) Hadamard block has nonzero determinant.  A
nonzero \(1\times1\) block has rank \(1\).  Rank is additive for block
diagonal matrices.
\end{proof}

\begin{proposition}[The mod-\(2\) Mukai pairing is nondegenerate]
\label{prop:qtet-mukai-mod2-nondegenerate}
Let \(\Lambda_{\mathrm{Muk}}\) be the even unimodular K3 Mukai lattice
with its integral bilinear form.  The reduction
\[
  \Lambda_{\mathrm{Muk}}/2\Lambda_{\mathrm{Muk}}
\]
with the induced \(\mathbb F_2\)-bilinear form is nondegenerate.
Consequently the Fourier matrix
\[
  S_{\lambda,\mu}
  =
  2^{-12}(-1)^{\langle\lambda,\mu\rangle_{\mathrm{Muk}}}
  \qquad
  \lambda,\mu\in\Lambda_{\mathrm{Muk}}/2\Lambda_{\mathrm{Muk}}
\]
has full rank \(2^{24}\).
\end{proposition}

\begin{proof}
Choose an integral Gram matrix \(A\) for the Mukai form.  Unimodularity
means \(\det A=\pm1\).  Reducing modulo \(2\) gives
\(\det(A\bmod 2)=1\), hence the bilinear form over \(\mathbb F_2\) is
nondegenerate.  The displayed matrix is the character table of the
finite abelian group \((\mathbb Z/2)^{24}\) under this perfect pairing;
character tables are invertible.
\end{proof}

\begin{corollary}[The \(324\) and \((\mathbb Z/2)^{24}\) pictures are
different]
\label{cor:qtet-324-vs-2-24}
A category with simple objects indexed by
\((\mathbb Z/2)^{24}\) has \(2^{24}\) simples.  A
\(324\times324\) matrix can describe a charge-\(2\) Fock truncation,
a \(324\)-object quotient, or a nonpointed category with \(324\)
simples.  The full pointed metric-group matrix for
\((\mathbb Z/2)^{24}\) has \(2^{24}\) labels.
\end{corollary}

\subsection{The root-of-unity statement}
\label{subsec:qtet-root-unity-statement}

\begin{theorem}[The \(N=2\) specialisation]
\label{thm:qtet-proved-n2}
At the point \((q_1,q_2,q_3)=(-1,-1,1)\):
\begin{enumerate}[label=\textup{(\roman*)}]
\item the Ding-Iohara current kernel in the displayed chart is
\(\omega_{\DIM}(x)=(1+x)^2/(1-x)\);
\item the positive-half identity
\(\CoHA(\mathbb C^3)\simeq Y^+(\gghone)\) remains a positive-half
statement; the full double requires Cartan and negative halves;
\item the K3 quantum-toroidal branch requires reduced equivariance;
toric \((\mathbb C^\times)^3\) weights enter through external charts;
\item \(324_{\mathrm{Hilb}}=324_{\mathrm{Fock}}=324\);
\item a Hadamard block decomposition of the stated type is full rank;
\item the full pointed Fourier matrix for
\((\mathbb Z/2)^{24}\) has rank \(2^{24}\).
\end{enumerate}
\end{theorem}

\begin{proof}
Parts (i), (iv), (v), and (vi) are the computations above.  Part (ii)
is Theorem~\ref{thm:qtet-coha-positive-half} and
Corollary~\ref{cor:qtet-positive-half-no-hopf}.  Part (iii) is
Remark~\ref{rem:qtet-k3-equivariance}.
\end{proof}

\begin{conjecture}[Root-of-unity simple modules]
\label{conj:qtet-root-unity-simples}
There exists a specified semisimple or logarithmic root-of-unity
category \(\mathcal C_{N=2}\), attached either to the K3
stable-envelope positive branch or to the completed BKM small form,
whose simple-object count is
\[
  |\Irr(\mathcal C_{N=2})|=324.
\]
The conjecture includes the construction of the functor from the
charge-\(2\) K3 Fock space to \(\mathcal C_{N=2}\) and the proof that
the images of the \(324\) Fock basis states are precisely the simple
objects after the chosen semisimplification or logarithmic quotient.
\end{conjecture}

\begin{conjecture}[Root-of-unity modular data]
\label{conj:qtet-root-unity-modular-data}
The modular data of \(\mathcal C_{N=2}\) is obtained from an explicitly
constructed root-of-unity \(R\)-matrix and its double-braiding trace.
If the simple set has cardinality \(324\), the \(S\)-matrix is a
\(324\times324\) matrix with labels, fusion rules, and rank supplied by
that category.  The full pointed metric-group data of
\((\mathbb Z/2)^{24}\) has \(2^{24}\) labels.  Any rank defect must be
computed from the \(324\)-object category itself.
\end{conjecture}

\subsection{Proof obligations}
\label{subsec:qtet-proof-obligations}

\begin{enumerate}[label=\textup{(\arabic*)}]
\item Construct the K3 quantum-toroidal algebra with its allowed
equivariance: generic K3, ADE/Kummer stratum, or \(K3\times E\) with
elliptic stable-envelope data.
\item Specify the root-of-unity algebra at \(N=2\): K3 positive
stable-envelope specialisation, BKM small form, or completed
Hall-Drinfeld double.  These constructions have distinct inputs and
distinct completions.
\item Give the integral form, quotient, and Cartan completion before
claiming central-rank behaviour at \(q=-1\).
\item Prove the passage from the \(324\)-dimensional charge-\(2\) Fock
space to \(324\) simple objects, including the semisimplification or
logarithmic quotient.
\item An \(S\)-matrix assertion must be one of two precise statements:
either the full pointed Fourier matrix on
\((\mathbb Z/2)^{24}\), which has \(2^{24}\) labels and full rank, or a
separately constructed \(324\times324\) modular matrix.
\item Keep the Mukai K3 Yangian branch separate from the
\(K3\times E\) BKM Hall-Drinfeld-double branch until the positive
boundary map and the completed double have both been constructed.
\end{enumerate}
